\documentclass[10pt]{article}
\usepackage[T1]{fontenc}
\usepackage[utf8]{inputenc}
\usepackage[french]{babel}
\usepackage{amsmath,amssymb,mathrsfs,mathpazo}
\usepackage{geometry}
\geometry{a4paper,margin=2cm}
\usepackage{graphicx,array,tabularx,tabularray,colortbl,multirow,multicol,booktabs,enumitem,pifont,ulem,fancybox,diagbox,cancel}
\usepackage{pstricks,pstricks-add,pst-plot,pst-tree,pst-node,pst-text,pst-func,pst-math,pst-bspline,pst-3dplot,pst-eucl,pst-solides3d}
\let\pscalculate\undefined
\usepackage{pst-bezier}
\usepackage{tikz,pgfplots,tkz-tab}
\usetikzlibrary{arrows,arrows.meta,patterns,calc,positioning,trees,decorations.markings,shapes.geometric,shapes.misc,shapes.arrows,intersections,decorations.pathmorphing,decorations.pathreplacing,backgrounds,shadows,babel}
\pgfplotsset{compat=1.18}
\usepgfplotslibrary{fillbetween}
\usepackage[np]{numprint}
\usepackage{listings,makecell,pas-tableur,twoopt}
\let\barre\undefined
\usepackage[pstricks]{tablvar}
\definecolor{darkgreen}{RGB}{0,100,0}
\providecommand{\red}{\color{red}}
\providecommand{\blue}{\color{blue}}
\providecommand{\green}{\color{green}}
\providecommand{\black}{\color{black}}
\providecommand{\white}{\color{white}}
\providecommand{\gray}{\color{gray}}
\providecommand{\cyan}{\color{cyan}}
\providecommand{\magenta}{\color{magenta}}
\providecommand{\yellow}{\color{yellow}}
\providecommand{\violet}{\color{violet}}
\providecommand{\orange}{\color{orange}}
\providecommand{\euro}{\texteuro}
\providecommand{\EUR}{\texteuro}
\providecommand{\ieme}{\textsuperscript{e}}
\providecommand{\degres}{\ensuremath{^\circ}}
\providecommand{\degree}{\ensuremath{^\circ}}
\providecommand{\vv}[1]{\overrightarrow{#1}}
\providecommand{\Coord}[1]{#1}
\providecommand{\decosix}{}
\providecommand{\starredbullet}{\ensuremath{\bullet}}
\pagestyle{empty}
\setlength{\parindent}{0pt}
\usepackage{adjustbox}\DeclareUnicodeCharacter{2236}{:}\DeclareUnicodeCharacter{2190}{\ensuremath{\leftarrow}}
\providecommand{\R}{\mathbb{R}}
\providecommand{\N}{\mathbb{N}}
\providecommand{\Z}{\mathbb{Z}}
\providecommand{\Q}{\mathbb{Q}}
\providecommand{\D}{\mathbb{D}}
\providecommand{\C}{\mathbb{C}}
\providecommand{\e}{{\mathrm{e}}}
\providecommand{\vect}[1]{\overrightarrow{#1}}
\providecommand{\vectt}[1]{\overrightarrow{\mathrm{#1}}}
\providecommand{\vvt}[1]{\overrightarrow{\mathrm{#1}}}
\providecommand{\np}[2][]{{#2~#1}}
\providecommand{\vv}[1]{{\overrightarrow{#1}}}
\renewcommand{\d}{{\mathrm{d}}}
\providecommand{\degres}{{^\circ}}
\providecommand{\degre}{{^\circ}}
\providecommand{\blue}{\color{blue}}
\providecommand{\red}{\color{red}}
\providecommand{\enskip}{\quad}
\providecommand{\llbracket}{\lbrack\!\lbrack}
\providecommand{\rrbracket}{\rbrack\!\rbrack}
\providecommand{\euro}{{\text{€}}}
\providecommand{\no}{{\text{n°}}}
\providecommand{\interval}[2]{{[#1;#2]}}
\providecommand{\intervFO}[2]{{[#1;#2[}}
\providecommand{\dotfill}{\ldots}

\providecommand{\ds}{\displaystyle}
\providecommand{\textding}[1]{#1}
\providecommand{\Oij}{\(\mathrm{O};\vec{\imath},\vec{\jmath}\)}
\providecommand{\Oijk}{\(\mathrm{O};\vec{\imath},\vec{\jmath},\vec{k}\)}
\providecommand{\Ouv}{\(\mathrm{O};\vec{u},\vec{v}\)}
\providecommand{\intervFF}[2]{\ensuremath{[#1;#2]}}
\renewcommand{\euro}{\ifmmode\text{\texteuro}\else\texteuro\fi}
\DeclareUnicodeCharacter{00B0}{\ifmmode^{\circ}\else\textdegree\fi}
\DeclareUnicodeCharacter{20AC}{\ifmmode\text{\texteuro}\else\texteuro\fi}
% Énoncés : Éducation nationale. Transcription et corrigés : APMEP, Denis Vergès et contributeurs. Découpage et classement : CoopMaths, Sébastien Lozano, Éric Elter et contributeurs. Adaptation de la présentation : Maths Explorer.
% CC BY-SA : https://coopmaths.fr/www/about/
% https://coopmaths.fr/alea/static/bac/2023/tex/bac_2023_09_sujet1_polynesie_2.tex
\begin{document}
\noindent\textit{APMEP / CoopMaths. CC BY-SA. Présentation adaptée par Maths Explorer.}\par\medskip
\medskip

\emph{Les parties A et B peuvent être traitées indépendamment}

\medskip
\textbf{Partie A}
\medskip

On considère la fonction $f$ définie sur $\R$ par

\[f(x) = \left(x + \dfrac12\right)\e^{-x} + x.\]

\medskip

\begin{enumerate}
\item Déterminer les limites de $f$ en $- \infty$ et en $+ \infty$.
\item On admet que $f$ est deux fois dérivable sur $\R$.
	\begin{enumerate}
		\item Démontrer que, pour tout $x \in \R$,
		
\[f''(x) = \left(x - \dfrac32\right) \e^{-x}.\]

		\item En déduire les variations et le minimum de la fonction $f'$ sur $\R$.
		\item Justifier que pour tout $x \in \R,\: f'(x) > 0$.
		\item En déduire que l'équation $f(x) = 0$ admet une unique solution sur $\R$.
		\item Donner une valeur arrondie à $10^{-3}$ de cette solution.
	\end{enumerate}
\end{enumerate}

\medskip
\textbf{Partie B}
\medskip

On considère une fonction $h$, définie et dérivable sur $\R$, ayant une expression de la forme 

\begin{center}$h(x)= (ax + b)\e^{-x} + x$, \:où $a$ et $b$\: sont deux réels.\end{center}

Dans un repère orthonormé ci-après figurent:

\begin{itemize}
\item[$\bullet$~~]la courbe représentative $\mathcal{C}_h$ de la fonction $h$ ;
\item[$\bullet$~~]les points A et B de coordonnées respectives $(-2~;~-2,5)$ et (2~;~3,5).
\end{itemize}

\begin{center}
\psset{unit=0.8cm}
\begin{pspicture*}(-5.9,-4.9)(5.9,5.9)
\psaxes[linewidth=1.25pt,labelFontSize=\scriptstyle](0,0)(-5.9,-4.9)(5.9,5.9)
\psplot[plotpoints=500,linewidth=1.25pt]{-4}{4}{1.5 x mul 0.5 add}
\psplot[plotpoints=2000,linewidth=1.25pt,linecolor=cyan]{-4}{4}{x 0.5 add 2.71828 x exp div x add}
\psdots(-2,-2.5)(2,3.5)
\uput[r](-2,-2.5){A}\uput[r](2,3.5){B}\uput[dr](3.8,4){\cyan $\mathcal{C}_h$}
\end{pspicture*}
\end{center}

\begin{enumerate}
\item Conjecturer, avec la précision permise par le graphique, les abscisses des éventuels points d'inflexion de la courbe représentative de la fonction $h$.
\item Sachant que la fonction $h$ admet sur $\R$ une dérivée seconde d'expression

\[h''(x) = - \dfrac32 \e^{- x} + x\e^{- x}.\]

valider ou non la conjecture précédente.
\item Déterminer une équation de la droite (AB).
\item Sachant que la droite (AB) est tangente à la courbe représentative de la fonction
$h$ au point d'abscisse 0, en déduire les valeurs de $a$ et $b$.
\end{enumerate}

\bigskip

\end{document}
